% =====================================================================
\lezione{09/10/2026}{Laboratorio 1 -- Severity assessment del caso Arvenza}
  {Prof.~Giuseppe Vaciago; Gianluca Gilardi (supporto)}
  {organizzazione dei gruppi; svolgimento del Lab~1 (data breach della mailbox paghe, metodologia
   ENISA); chiarimenti sul caso durante il lavoro}
\label{lez:2026-10-09}
% =====================================================================

\begin{nota}[Fonte]
Trascrizione automatica della registrazione: la sessione è in gran parte silenziosa (gruppi al
lavoro). Sono riportate solo le indicazioni udibili del docente. Il caso, la metodologia e la
soluzione sono trattati nel Capitolo~\ref{lab:01}.
\end{nota}

\section{Sintesi}
Primo laboratorio, gestito su Zoom (il docente era collegato da remoto). Dopo la registrazione dei gruppi (alla fine circa 25), gli
studenti, nel ruolo di DPO di Arvenza Industrie, preparano in circa un'ora e mezza un
\textbf{severity assessment} con metodologia ENISA e un parere che concili le posizioni di CEO,
IT Manager e HR Director. Il docente e Gianluca rispondono alle domande.

\section{Organizzazione}
\begin{itemize}
  \item Ogni gruppo comunica in chat nome del gruppo e componenti.
  \item \textbf{Consegna}: un solo componente invia l'email al docente, con gli altri in copia; nel
    documento nomi e \textbf{numero di matricola} di tutti. Domande ammesse fino alle 12:45--12:50,
    consegna entro le \textbf{13:15}; oltre quell'ora il lavoro non è considerato.
  \item Il docente sperimenta una sala secondaria su Zoom per le domande, ma conclude che è meglio
    avere tutte le domande nella sala principale, perché le risposte servono a tutti.
  \item La valutazione arriva entro il weekend; i gruppi migliori presentano lunedì~12/10 (cinque
    minuti circa ciascuno, in tutto una quarantina di minuti).
\end{itemize}

\section{Indicazioni del docente durante il lab}
\begin{itemize}
  \item L'obiettivo è il severity assessment ``come presentato lunedì'', combinando le esigenze dei
    tre punti di vista. \textbf{Non c'è una regola fissa}: la scelta giusta è quella coerente con il
    principio di \emph{accountability}.
  \item Chiarimenti forniti a tutti sui dati del caso:
    \begin{itemize}
      \item i dati bancari (anche quelli dei tre dipendenti impersonati) \textbf{non erano cifrati};
      \item le buste paga 2023 vanno lette con attenzione perché la loro situazione è diversa da
        quella delle 2024--26 (le password di gennaio--aprile erano nella stessa mailbox);
      \item i documenti sanitari e i CV \textbf{non erano cifrati}.
    \end{itemize}
\end{itemize}

\begin{esame}
Il ragionamento conta più dell'opzione scelta: ogni punteggio deve poggiare su un fatto, e ogni
scostamento dal valore di base va motivato. Lo svolgimento completo è nel Capitolo~\ref{lab:01}.
\end{esame}

\section{Domande aperte e punti da verificare}
\begin{verifica}
\begin{itemize}
  \item Il testo del caso annuncia una \textbf{Parte~B} (aggiornamento dell'8 ottobre) dopo il
    debrief: non risulta né dalla registrazione né dal materiale nella cartella del lab.
  \item Orari esatti di chiusura delle domande e di consegna: ricostruiti dal parlato (12:45--12:50 e
    13:15).
\end{itemize}
\end{verifica}
